\chapter{Superando o mundo fechado}
\label{cap:closedworld}

\begin{objetivos}
Ao final deste capítulo, você será capaz de:
\begin{itemize}[leftmargin=*,itemsep=0.2em]
      \item descrever e interpretar um experimento de mundo aberto por classe escondida;
  \item analisar o compromisso entre comprimento de janela, acurácia e tempo de resposta
        operacional;
  \item interpretar a validação de um classificador contra um modelo físico independente
        e analisar as discrepâncias;
  \item estimar o custo computacional de um ensemble e discutir sua viabilidade no
        ambiente avaliado;
  \item posicionar corretamente as limitações de um resultado positivo.
\end{itemize}
\end{objetivos}

\section{Do princípio à evidência}

O \cref{cap:binarios} estabeleceu que um ensemble de classificadores binários
\emph{pode} rejeitar, sem perda observada na acurácia agregada daquele experimento.
Faltava confrontá-lo com uma classe que o sistema realmente nunca viu e medir se a
rejeição continuava útil nesse protocolo.

Este capítulo apresenta essa evidência para uma classe escondida, junto com três investigações que a
acompanham: a varredura sistemática do comprimento de janela, a validação contra um
modelo termodinâmico independente e a análise do custo computacional. É o
aprofundamento do \cref{cap:binarios} --- mesmos classificadores, perguntas mais duras.

\section{O experimento da classe escondida}

\subsection{Protocolo}

O experimento é conceitualmente simples e metodologicamente decisivo.

\begin{enumerate}[leftmargin=*,itemsep=0.3em]
  \item Treinam-se classificadores binários para \textbf{seis} classes apenas.
  \item A classe 8 --- hidrato na linha de produção --- é \textbf{completamente removida}
        do treinamento. Nenhuma janela dela é vista, em nenhum momento, por nenhum
        classificador.
  \item Apresentam-se ao ensemble as instâncias da classe 8.
  \item Mede-se a fração dessas instâncias para as quais \emph{todos} os seis escores
        ficam abaixo do limiar.
\end{enumerate}

A pergunta que o experimento responde não é ``o sistema classifica bem?'', mas
\textbf{``o sistema sabe que não sabe?''}.

\subsection{Resultado}

\begin{resultadochave}
$\bm{89\,\%}$ das instâncias da classe escondida foram sinalizadas como ``nenhuma das
classes treinadas''.
\end{resultadochave}

A \cref{fig:closed-multiclasse} mostra o que acontece quando o mesmo experimento é feito
com um classificador multiclasse, e é a ilustração mais direta do problema que este
livro ataca.

\begin{figure}[htbp]
\centering
\includegraphics[width=0.9\textwidth]{Figure_7.png}
\caption{Escores de um classificador multiclasse treinado em seis classes, quando
apresentado às instâncias da classe 8 --- que ele nunca viu. Note a concentração de
massa próxima de $1{,}0$ para as classes 3, 2 e 5: o modelo não hesita, ele afirma com
convicção máxima que o hidrato é outra coisa. A restrição
$\sum_k \hat{p}_k = 1$ não lhe permite outra resposta. Adaptado de
\citet{Lopes2026PST}.}
\label{fig:closed-multiclasse}
\end{figure}

\begin{destaque}[O que a figura mostra e o que ela não mostra]
Ela não mostra um modelo mal treinado. Os picos próximos de $1{,}0$ indicam alta
confiança, e a alta confiança é apropriada \emph{dentro} do mundo fechado em que o
modelo foi treinado e avaliado.

O que a figura mostra é que a confiança de um classificador multiclasse não carrega
informação alguma sobre se a entrada pertence ao domínio para o qual ele foi treinado.
É um alerta que vale muito além do monitoramento de poços.
\end{destaque}

\begin{destaque}[Por que este é o resultado central da tese]
Um classificador multiclasse, no mesmo experimento, teria acertado $0\,\%$ --- não por
ser mal treinado, mas por construção. A restrição $\sum_k \hat{p}_k = 1$ o obriga a
distribuir toda a massa de probabilidade entre as seis classes que conhece. A resposta
``nenhuma delas'' não pertence ao seu espaço de saída.

Os $89\,\%$ não medem, portanto, uma melhoria incremental sobre uma linha de base: eles
medem a existência de uma capacidade que antes era literalmente inexpressável.
\end{destaque}

\subsection{Os onze por cento restantes}

Um resultado honesto exige olhar para o complemento. Onze por cento das instâncias de
hidrato foram atribuídas com confiança a alguma das seis classes conhecidas --- ou seja,
o sistema afirmou algo falso com convicção.

Em termos operacionais, esse é o modo de falha mais perigoso do sistema. Um alarme
ausente leva o operador a investigar por conta própria; um diagnóstico errado o leva a
investigar \emph{na direção errada}. O \cref{cap:operacao} retornará a essa distinção ao
discutir o projeto da interface.

\begin{destaque}[Uma comparação de escala]
Vale registrar a ordem de grandeza. Os $89\,\%$ deste experimento contrastam com os
$52\,\%$ obtidos sobre a anomalia de ICV no \cref{cap:mundoaberto} --- e a diferença é
instrutiva.

O hidrato tem assinatura física própria e distintiva. A falha de ICV é sintomaticamente
próxima de classes existentes. A capacidade de rejeição, portanto, não é uma propriedade
fixa do método: ela depende de quão distinta é a novidade em relação ao que já se
conhece. Sistemas de mundo aberto detectam melhor o que é mais estranho --- o que é
tranquilizador e, ao mesmo tempo, é a razão pela qual o número reportado em qualquer
trabalho da área deve vir acompanhado da identificação da classe escondida.
\end{destaque}

\section{O comprimento da janela}

\subsection{Um compromisso que não é técnico}

Quanto tempo de sinal é necessário para reconhecer um hidrato? Em uma janela causal,
$T$ define o histórico mínimo necessário após o início da aquisição; o atraso adicional
depende do passo entre janelas, do alinhamento do evento e da frequência de inferência.

A \cref{tab:closed-janelas} apresenta a varredura completa, de \SI{150}{\second} a
\SI{7500}{\second}, para o classificador binário de hidrato.

\begin{table}[htbp]
\centering
\caption{Desempenho do classificador binário de hidrato em função do comprimento da
janela $T$. FP e FN são as taxas de falso positivo e falso negativo.}
\label{tab:closed-janelas}
\small
\begin{tabular}{rcccccc}
\toprule
$\bm{T}$ \textbf{(s)} & \textbf{ACC} & \textbf{PR} & \textbf{REC} & $\bm{F_1}$ &
\textbf{FP} & \textbf{FN} \\
\midrule
150   & $0{,}866$ & $0{,}863$ & $0{,}898$ & $0{,}881$ & $0{,}175$ & $0{,}102$ \\
300   & $0{,}891$ & $0{,}883$ & $0{,}928$ & $0{,}905$ & $0{,}154$ & $0{,}072$ \\
600   & $0{,}891$ & $0{,}891$ & $0{,}918$ & $0{,}904$ & $0{,}144$ & $0{,}082$ \\
900   & $0{,}902$ & $0{,}885$ & $0{,}952$ & $0{,}917$ & $0{,}166$ & $0{,}048$ \\
1200  & $0{,}898$ & $0{,}893$ & $0{,}937$ & $0{,}915$ & $0{,}157$ & $0{,}063$ \\
1500  & $0{,}907$ & $0{,}911$ & $0{,}935$ & $0{,}923$ & $0{,}132$ & $0{,}065$ \\
2250  & $0{,}899$ & $0{,}923$ & $0{,}910$ & $0{,}917$ & $0{,}117$ & $0{,}090$ \\
3000  & $0{,}930$ & $0{,}949$ & $0{,}941$ & $0{,}945$ & $0{,}090$ & $0{,}059$ \\
3750  & $0{,}961$ & $0{,}958$ & $0{,}985$ & $0{,}971$ & $0{,}085$ & $0{,}015$ \\
4500  & $0{,}920$ & $0{,}895$ & $0{,}997$ & $0{,}943$ & $0{,}236$ & $\bm{0{,}003}$ \\
5250  & $0{,}962$ & $\bm{0{,}977}$ & $0{,}966$ & $0{,}971$ & $\bm{0{,}047}$ & $0{,}034$ \\
6000  & $0{,}869$ & $0{,}968$ & $0{,}832$ & $0{,}895$ & $0{,}056$ & $0{,}168$ \\
6750  & $0{,}942$ & $0{,}922$ & $\bm{1{,}000}$ & $0{,}959$ & $0{,}180$ & $\bm{0{,}000}$ \\
7500  & $\bm{0{,}966}$ & $0{,}958$ & $0{,}994$ & $\bm{0{,}975}$ & $0{,}092$ & $0{,}006$ \\
\bottomrule
\end{tabular}
\end{table}

\subsection{Três leituras da tabela}

\paragraph{A tendência.} A acurácia cresce de $0{,}866$ para $0{,}966$ ao longo de duas
ordens de grandeza de $T$. O ganho é real, mas modesto: dez pontos percentuais para um
aumento de cinquenta vezes no comprimento da janela. E ele não é monótono --- há uma
queda pronunciada em \SI{6000}{\second} ($0{,}869$), com revocação despencando para
$0{,}832$.

\begin{destaque}[Sobre a não monotonicidade]
A queda em \SI{6000}{\second} não deve ser explicada apressadamente. Duas causas
concorrem. Primeira, janelas mais longas produzem \emph{menos} exemplos de treinamento a
partir do mesmo conjunto de dados --- a variância da estimativa cresce. Segunda, com
$T = \SI{6000}{\second}$ a janela passa a conter, frequentemente, tanto o regime normal
quanto o evento inteiro, e o rótulo agregado torna-se ambíguo.

A lição metodológica: uma varredura de hiperparâmetro que produz curva ruidosa está
comunicando algo sobre a quantidade de dados, não apenas sobre o hiperparâmetro.
\end{destaque}

\paragraph{O compromisso entre FP e FN.} As duas últimas colunas revelam um padrão que a
acurácia esconde. Em \SI{4500}{\second}, a taxa de falso negativo cai para $0{,}003$ ---
praticamente nenhum hidrato escapa --- mas a de falso positivo sobe para $0{,}236$. Em
\SI{5250}{\second}, quase o oposto: FP de $0{,}047$ e FN de $0{,}034$.

Qual configuração é melhor? \textbf{A pergunta não tem resposta técnica.} Um hidrato não
detectado pode custar dias de produção e uma intervenção submarina; um alarme falso custa
atenção do operador. A razão entre esses custos é uma decisão da operadora, e o papel do
engenheiro é apresentar a curva, não escolher o ponto.

\paragraph{A leitura operacional.} É a mais importante, e a \cref{tab:closed-tempo} a
torna explícita.

\begin{table}[htbp]
\centering
\caption{Compromisso entre histórico mínimo e acurácia para três comprimentos de janela
representativos; o atraso efetivo depende do passo e do alinhamento.}
\label{tab:closed-tempo}
\begin{tabular}{rcc}
\toprule
$\bm{T}$ \textbf{(s)} & \textbf{Histórico mínimo} & \textbf{Acurácia} \\
\midrule
150  & $2{,}5$ minutos & $86{,}6\,\%$ \\
300  & $5$ minutos     & $89{,}1\,\%$ \\
900  & $15$ minutos    & $90{,}2\,\%$ \\
\bottomrule
\end{tabular}
\end{table}

\begin{resultadochave}
Reduzir a janela de \SI{900}{\second} para \SI{150}{\second} --- de quinze para dois
minutos e meio de histórico --- custa $3{,}6$ pontos percentuais de acurácia neste teste.
\end{resultadochave}

Esse número precisa ser confrontado, em cada implantação, com o intervalo disponível
para agir. A faixa de trinta minutos a cinco horas antes atribuída a
\citet{VARGAS2019} não foi confirmada na fonte consultada nesta revisão e, por isso, não
é usada aqui como evidência.

\begin{destaque}[A conclusão que inverte a intuição]
O gargalo do sistema não é o tempo de detecção. É a confiabilidade do diagnóstico.

Otimizar por latência, neste problema, é otimizar a variável errada. É por isso que o
restante deste livro investe em segmentação precisa (\cref{cap:segmentacao}), em
detecção de novidade confiável (\cref{cap:mahalanobis}) e em explicabilidade
(\cref{parte:explicabilidade}) --- e não em inferência mais rápida.
\end{destaque}

\section{Validação contra um modelo físico}

\subsection{A ideia}

Comparar um classificador com outro classificador diz pouco: ambos podem estar
compartilhando os mesmos vieses do conjunto de dados. Comparar um classificador com um
\textbf{modelo termodinâmico independente} é outra coisa --- o modelo físico não foi
treinado em dado nenhum e não conhece o 3W.

\subsection{O modelo de Motiee}

\begin{figure}[htbp]
\centering
\includegraphics[width=\textwidth]{Figure_9.png}
\caption{Arquitetura de validação cruzada entre modelo físico e modelo aprendido. Os
mesmos valores de sensor alimentam dois caminhos independentes --- a análise
pressão--temperatura contra a envoltória de formação de hidrato e o classificador
binário --- cujos resultados são então confrontados antes de gerar o alerta. Adaptado de
\citet{Lopes2026PST}.}
\label{fig:closed-arquitetura}
\end{figure}

A correlação empírica de \citet[p.~98--99]{montie1991} estima a temperatura de formação
de hidrato $T_{\text{hid}}$ a partir da pressão absoluta $P$ e da densidade relativa do
gás $q$. Na forma publicada, define-se $x=\log_{10}(P/\mathrm{psia})$, e a saída é dada
em graus Fahrenheit:
\begin{equation}
\begin{split}
      T_{\text{hid},\,{}^\circ\mathrm{F}} = {} & -238{,}24469 + 78{,}99667\,x
            - 5{,}352544\,x^2 \\
      & + 349{,}473877\,q - 150{,}854675\,q^2 - 27{,}604065\,xq.
\end{split}
\label{eq:motiee}
\end{equation}
Trata-se de uma correlação de triagem para gás natural doce baseada em gravidade
específica, não de um modelo composicional. Por isso, seu uso fora da faixa de pressão e
composição coberta pela fonte exige validação contra dados de equilíbrio ou simulador
termodinâmico.

A partir dela define-se um índice de proximidade da condição de hidrato,
\begin{equation}
  C_{\text{hid}} = 100 \, \frac{T - T_{\text{hid}}}{T},
  \label{eq:chid}
\end{equation}
onde $T$ é a temperatura medida. Valores negativos ou próximos de zero indicam operação
dentro ou junto da envoltória de formação. A curva pressão--temperatura de referência é a
de \citet{Safamirzaei2015}.

\begin{figure}[htbp]
\centering
\includegraphics[width=0.92\textwidth]{figure_5.eps}
\caption{Envoltória de formação de hidrato no plano temperatura--pressão. A curva separa
a zona de operação normal da zona de formação: para uma dada pressão, temperaturas à
esquerda da curva colocam o sistema em risco. Adaptado de \citet{Lopes2026PST}.}
\label{fig:closed-envoltoria}
\end{figure}

A \cref{fig:closed-envoltoria} apresenta essa envoltória. Note que a curva é monotônica e
suave --- e é justamente essa simplicidade que a torna um instrumento de validação útil.
Ela não tem parâmetros ajustados aos dados, não conhece o 3W e não pode ter aprendido
nenhum atalho do conjunto. Se o classificador concorda com ela, a concordância significa
alguma coisa.

\begin{exemplo}{Usando a \cref{eq:motiee} na prática}
Considere um trecho de linha operando a $P = \SI{7}{\mega\pascal}$, ou aproximadamente
$1015\,\text{psia}$, com gás de densidade relativa $q = 0{,}65$ e temperatura medida
$T = \SI{18}{\celsius}$.

Com $x=\log_{10}(1015)\approx3{,}007$, a \cref{eq:motiee} fornece
$T_{\text{hid}}\approx60{,}4\,{}^\circ\mathrm{F}\approx\SI{15,8}{\celsius}$. Pela
\cref{eq:chid}, $C_{\text{hid}} \approx 100 \times (18 - 15{,}8)/18 \approx 12\,\%$.

Uma queda de aproximadamente $2{,}2$ graus na temperatura da linha levaria
$C_{\text{hid}}$ a zero. É esse tipo de raciocínio de margem que o modelo analítico
oferece e que um classificador, por si só, não comunica.
\end{exemplo}

\subsection{Resultado da comparação}

\begin{table}[htbp]
\centering
\caption{Classificador binário de hidrato \emph{versus} critério analítico baseado na
correlação de Motiee.}
\label{tab:closed-analitico}
\begin{tabular}{lcc}
\toprule
\textbf{Métrica} & \textbf{Classificador binário} & \textbf{Critério analítico} \\
\midrule
Acurácia            & $\bm{90{,}02\,\%}$ & $76{,}60\,\%$ \\
Falsos positivos    & $16{,}80\,\%$      & $18{,}61\,\%$ \\
Falsos negativos    & $\bm{4{,}80\,\%}$  & $20{,}10\,\%$ \\
\bottomrule
\end{tabular}
\end{table}

\begin{resultadochave}
O classificador supera o critério analítico em treze pontos percentuais de acurácia. A
diferença concentra-se nos falsos negativos: $4{,}8\,\%$ contra $20{,}1\,\%$ --- uma
redução de quatro vezes.
\end{resultadochave}

\subsection{Interpretando a diferença}

Por que o modelo físico perde? Não por estar errado. A correlação de Motiee descreve
corretamente a \emph{condição termodinâmica de formação} de hidrato. O que ela não
descreve é a \emph{cinética}: a formação leva tempo, depende da presença de água livre,
de nucleação, de histórico de escoamento e de eventual injeção de inibidor.

Um trecho pode estar dentro da envoltória por horas sem formar hidrato; outro pode formar
com margem aparentemente confortável, se houver acúmulo de água. O modelo analítico
responde ``é termodinamicamente possível?''; o classificador aprende a responder ``está
acontecendo?''.

\begin{destaque}[Complementaridade, não substituição]
As taxas de falso positivo são quase idênticas --- $16{,}8\,\%$ e $18{,}6\,\%$ --- o que
sugere que ambos os métodos são igualmente conservadores ao alertar. A diferença está
inteiramente na detecção.

A arquitetura sensata não escolhe entre os dois. O modelo analítico fornece contexto
físico interpretável --- a margem de $C_{\text{hid}}$ --- que o classificador não produz;
o classificador fornece a sensibilidade que o modelo não tem. O
\cref{cap:agente} retomará essa ideia em outro registro, ao alimentar o agente de
linguagem com conhecimento de domínio explícito.
\end{destaque}

Vale ainda notar que esta validação é possível apenas para o hidrato, porque apenas ele
dispõe de um modelo físico fechado e bem estabelecido. Para golfadas, incrustação ou
falha de ICV não há equivalente --- e essa é uma limitação da estratégia de validação,
não do classificador.

\subsection{A complementaridade na tela do operador}

\begin{figure}[htbp]
\centering
\includegraphics[width=\textwidth]{Figure_10.png}
\caption{Painel operacional com os dois indicadores lado a lado. À direita, a curva de
formação de hidrato com o ponto de operação atual; abaixo, os dois medidores de risco:
o analítico (PVT), em $19{,}71\,\%$, e o aprendido (IA), em $0{,}00\,\%$. Os dois números
respondem a perguntas diferentes e são apresentados sem tentativa de fundi-los em um só.
Adaptado de \citet{Lopes2026PST}.}
\label{fig:closed-painel}
\end{figure}

A \cref{fig:closed-painel} materializa a recomendação da seção anterior. A decisão de
projeto --- exibir os dois riscos separadamente, em vez de combiná-los em um índice
único --- é deliberada. Um risco analítico de $19{,}71\,\%$ com risco aprendido nulo
comunica ao operador algo específico e acionável: \emph{as condições termodinâmicas se
aproximam da envoltória, mas ainda não há assinatura de formação}. Um índice combinado
de, digamos, $10\,\%$ destruiria essa informação.

\begin{destaque}[Fundir indicadores é perder informação]
A tentação de reduzir múltiplos indicadores a um número único é forte, e quase sempre
equivocada em contextos de decisão especializada. Dois indicadores que discordam são
informativos precisamente por discordarem. O \cref{cap:operacao} desenvolverá esse
princípio ao discutir o projeto do painel do operador.
\end{destaque}

\section{Custo computacional}

O argumento recorrente contra ensembles é o custo. Convém quantificá-lo.

\begin{table}[htbp]
\centering
\caption{Custo computacional por inferência, para janela de \SI{900}{\second}.}
\label{tab:closed-custo}
\begin{tabular}{lcc}
\toprule
& \textbf{Memória} & \textbf{Latência} \\
\midrule
Um classificador binário   & $\approx \SI{2,3}{\mega\byte}$ & $\approx \SI{45}{\milli\second}$ \\
Ensemble de nove           & $\approx \SI{21}{\mega\byte}$  & $\approx \SI{405}{\milli\second}$ \\
\bottomrule
\end{tabular}
\end{table}

\begin{resultadochave}
O ensemble completo ocupa cerca de \SI{21}{\mega\byte} e responde em cerca de
\SI{405}{\milli\second} para uma janela de \SI{900}{\second}.
\end{resultadochave}

A inferência é mais de três ordens de grandeza menor que a janela avaliada. Cabe em um computador industrial
embarcado, e o custo de multiplicar por nove --- que parecia uma objeção séria no
\cref{cap:binarios} --- não é limitante no hardware e na frequência avaliados.

\begin{destaque}[Uma observação sobre argumentos de eficiência]
Este é um caso em que a quantificação delimita o debate. Objeções de custo computacional
devem ser sempre confrontadas com a escala de tempo do processo monitorado. Em
monitoramento de hidrato no cenário avaliado, a inferência é rápida em relação à janela;
rede, aquisição, concorrência entre poços e requisitos de disponibilidade ainda precisam
ser medidos na implantação.
\end{destaque}

\section{Limitações}

\begin{enumerate}[leftmargin=*,itemsep=0.4em]

\item \textbf{Uma única classe escondida.} O experimento foi feito removendo a classe 8.
Não se sabe se $89\,\%$ se sustentaria com outra classe --- e a comparação com os
$52\,\%$ da ICV no \cref{cap:mundoaberto} sugere fortemente que não.

\item \textbf{Detecção sem localização.} O sistema classifica a janela inteira. Não diz
onde, dentro dela, o evento começa. É a lacuna que o \cref{cap:segmentacao} preenche.

\item \textbf{Rejeição sem organização.} Rejeitar é o primeiro passo; o
\cref{cap:mundoaberto} mostrou o que vem depois, e mostrou também que essa etapa é a
mais frágil da cadeia.

\item \textbf{Rejeição sem explicação.} O sistema diz ``desconhecido'' sem dizer por quê.
É a \cref{parte:explicabilidade}.

\item \textbf{Validação física restrita ao hidrato.} As demais classes não dispõem de
modelo independente contra o qual comparar.

\end{enumerate}

\resumocapitulo{%
O experimento decisivo do princípio de mundo aberto consiste em treinar o ensemble sobre
seis classes, esconder completamente a classe 8 --- hidrato na produção --- e verificar
se suas instâncias são reconhecidas como desconhecidas: $89\,\%$ delas são, contra
$0\,\%$ que um classificador multiclasse obteria por construção. A varredura do
comprimento de janela, de \SI{150}{\second} a \SI{7500}{\second}, mostra ganho modesto e
não monótono ($0{,}866$ a $0{,}966$) e revela que reduzir a janela de quinze para dois
minutos e meio de histórico custa $3{,}6$ pontos percentuais; o atraso operacional ainda
depende do passo e do alinhamento. No ambiente avaliado, a inferência não é o gargalo.
A validação contra a correlação
termodinâmica de Motiee dá ao classificador $90{,}02\,\%$ de acurácia contra $76{,}60\,\%$
do critério analítico, com falsos negativos de $4{,}8\,\%$ contra $20{,}1\,\%$: o modelo
físico descreve a condição de formação, o classificador aprende a ocorrência. O custo do
ensemble --- \SI{21}{\mega\byte} e \SI{405}{\milli\second} --- é três ordens de grandeza
menor que a escala de tempo assumida para o processo e não foi limitante no ambiente
avaliado.}

\begin{exercicios}
\item Repita o experimento da classe escondida para cada uma das nove classes,
      reportando a taxa de rejeição em cada caso. Correlacione os resultados com alguma
      medida de distinção da assinatura --- por exemplo, a distância entre centroides no
      espaço latente. A hipótese de que ``a rejeição depende de quão estranha é a
      novidade'' se sustenta?

\item A \cref{tab:closed-janelas} mostra queda de acurácia em \SI{6000}{\second}.
      Projete dois experimentos que discriminem entre as duas explicações oferecidas
      (menos exemplos de treinamento \emph{versus} rótulo agregado ambíguo).

\item Para $T = \SI{4500}{\second}$, FP $= 0{,}236$ e FN $= 0{,}003$; para
      $T = \SI{5250}{\second}$, FP $= 0{,}047$ e FN $= 0{,}034$. Formule o problema de
      escolha como minimização de custo esperado, atribua custos plausíveis a um hidrato
      não detectado e a um alarme falso, e determine para que razão de custos cada
      configuração é preferível.

\item \textbf{Projeto de reprodução.} Com acesso às instâncias e às variáveis físicas
      necessárias, implemente a \cref{eq:motiee} e a \cref{eq:chid} para as classes 8 e
      9 do 3W e busque reproduzir a \cref{tab:closed-analitico}. Registre protocolo e
      divergências. Em seguida, examine os casos em que o critério analítico erra e o
      classificador acerta: há padrão?

\item O modelo analítico e o classificador têm taxas de falso positivo semelhantes mas
      taxas de falso negativo muito distintas. Proponha três formas de combiná-los ---
      conjunção, disjunção e voto ponderado --- e avalie cada uma quanto a FP, FN e
      interpretabilidade.

\item Estime o custo computacional do ensemble para os duzentos e cinquenta poços
      mencionados no \cref{cap:dual}, supondo inferência a cada \SI{60}{\second}. Discuta
      arquiteturas de implantação: inferência na plataforma, em terra, ou híbrida.

\item Argumente a favor e contra a seguinte afirmação: ``o resultado de $89\,\%$ é
      superestimado porque a classe 8 tem assinatura particularmente distintiva''. Que
      evidência adicional resolveria a disputa?
\end{exercicios}

\begin{leituras}
\item \citet{Lopes2026PST} --- o artigo que originou este capítulo.
\item \citet{LopesTese2025} --- a tese, para o contexto completo.
\item \citet{montie1991} e \citet{Safamirzaei2015} --- as correlações termodinâmicas
      usadas na validação independente.
\item \citet{VARGAS2019} --- a descrição do conjunto 3W.
\item \citet{Scheirer2013} --- a formalização do risco de espaço aberto, que dá suporte
      teórico ao experimento da classe escondida.
\end{leituras}
